\section{The theorem and the precise foundational target}
\label{sec:foundations}

\begin{theorem}[Witnessed four-square theorem in $\ipra$]
\label{thm:main}
There are four explicitly specified primitive recursive functions
$W_i:\NN\to\NN$, $0\le i<4$, such that $\ipra$ proves
\begin{equation}
\label{eq:main}
 \forall n\;\bigl(W_0(n)^2+W_1(n)^2+W_2(n)^2+W_3(n)^2=n\bigr).
\end{equation}
It also proves $W_i(n)\le n$ for each $i$. Consequently it proves
\begin{equation}
\label{eq:exists}
 \forall n\,\exists a,b,c,d\le n\quad n=a^2+b^2+c^2+d^2.
\end{equation}
The convention is $\NN=\{0,1,2,\ldots\}$.
\end{theorem}

The functions will be defined before their correctness is concluded. Thus
the assertion is stronger than a pointwise existence argument: the witnesses
are supplied by fixed arithmetic programs, and the proof establishes their
correctness uniformly inside the specified theory.

\subsection{What is meant by intuitionistic PRA}

We use the standard first-order presentation of $\ipra$. Its only sort is
$\NN$. Its language has zero, successor, and a function symbol for every
primitive recursive definition, generated from projections by composition
and primitive recursion. If
\[
 f(0,\vec x)=g(\vec x),\qquad
 f(t+1,\vec x)=h(t,f(t,\vec x),\vec x),
\]
the corresponding defining equations are axioms. The non-logical axioms
also state that successor is injective and never zero, and include induction
\[
 \bigl(A(0,\vec u)\land
       \forall t(A(t,\vec u)\to A(t+1,\vec u))\bigr)
 \to \forall t\,A(t,\vec u)
\]
only when $A$ is quantifier-free. Free parameters are allowed.

The logic is intuitionistic first-order logic with equality: the usual
introduction and elimination rules for $\land,\lor,\to,\forall,\exists$,
falsehood elimination, and equality reflexivity and substitution. No general
law of excluded middle, choice principle, or double-negation elimination is
added. This is the presentation implemented by ProveIt's
\code{ConciseIPRA.PR} and \code{PRAAx}~\cite{proveit}.

Writing a new primitive recursive function in this presentation adds a
name for a definition, not a stronger existence axiom. A smaller fixed
language may instead compile the definitions into its established primitive
recursive term language.

\subsection{Decidable arithmetic and bounded statements}

Truncated subtraction $a\monus b=\max(a-b,0)$, equality tests, order tests,
parity, natural quotient and remainder, bounded sums, and bounded products
are primitive recursive. Their elementary laws have quantifier-free
induction proofs. The details of quotient/remainder and coding used here
are given in Appendix~\ref{app:coding}.

Every bounded predicate occurring in this report has a specified primitive
recursive characteristic bit. For a bit-valued primitive recursive
predicate $b(i,\vec u)$, define
\[
 \operatorname{All}(0,\vec u)=1,\qquad
 \operatorname{All}(k+1,\vec u)
   =\operatorname{All}(k,\vec u)b(k,\vec u).
\]
The equation $\operatorname{All}(t,\vec u)=1$ expresses the finite conjunction
$\bigwedge_{i<t}b(i,\vec u)=1$. A similar fold expresses finite disjunction.
Negation of a bit is $1\monus b$; conjunction is multiplication; disjunction
can be written $\min(1,b+c)$. Thus the finite tests are decidable in $\ipra$.

When a proof below inducts on a statement displayed with bounded
quantifiers, the actual induction formula is its characteristic-bit
equation. This is a concrete elimination of bounded quantifiers, not an
extension of the induction scheme to arbitrary quantified formulas.
Appendix~\ref{app:folds} gives the needed fold and search reasoning.

In particular, eliminating a failed bounded search is constructive. A
computed search flag is either $0$ or $1$. If the $0$ branch entails a
contradiction, the $1$ branch remains and contains actual recorded witnesses.
This is a decidable case split, not an application of double-negation
elimination to an arbitrary existential assertion.

\subsection{Signed integers and finite records are notation}

For readability we calculate with $\ZZ$ and four-tuples. Internally, a signed
integer is a natural code, and a tuple is a nested primitive recursive pair.
All operations used below, including signed Euclidean remainder and exact
division by a positive natural, have primitive recursive implementations.
Appendix~\ref{app:coding} gives an explicit encoding and operations.

For $x=(x_0,x_1,x_2,x_3)\in\ZZ^4$, put
\[
 \Nrm{x}=x_0^2+x_1^2+x_2^2+x_3^2.
\]
This is a natural-valued primitive recursive function of the tuple code.
The state assertion $\Nrm{x}=mp$ is therefore quantifier-free after coding.
No quantification over arbitrary functions, sets of integers, or infinite
sequences is needed anywhere in the proof.

\section{The algebraic identity that carries the construction}
\label{sec:identity}

\begin{definition}[Euler product]
For $x=(a,b,c,d)$ and $y=(e,f,g,h)$, define
\begin{equation}
\label{eq:product}
 x*y=\bigl(
 ae-bf-cg-dh,\;
 af+be+ch-dg,\;
 ag-bh+ce+df,\;
 ah+bg-cf+de
 \bigr).
\end{equation}
Write $\bar y=(e,-f,-g,-h)$.
\end{definition}

Only this explicit polynomial operation will be used. No theory of
quaternion algebras is required.

\begin{lemma}[Euler's four-square identity]
\label{lem:euler}
For all integer four-tuples $x,y$,
\[
 \Nrm{x*y}=\Nrm{x}\Nrm{y},\qquad
 x*\bar x=(\Nrm{x},0,0,0),\qquad
 \Nrm{\bar x}=\Nrm{x}.
\]
These are finite polynomial identities and have $\ipra$ proofs.
\end{lemma}

\begin{proof}
The four coordinates of $x*y$ are the rows of the following matrix applied
to $(e,f,g,h)^{\mathsf T}$:
\[
 L_x=\begin{pmatrix}
 a&-b&-c&-d\\
 b&a&-d&c\\
 c&d&a&-b\\
 d&-c&b&a
 \end{pmatrix}.
\]
Every column has squared length $a^2+b^2+c^2+d^2$. The six distinct
column dot products are
\begin{align*}
 -ab+ab+cd-dc&=0,&
 -ac-bd+ac+bd&=0,&
 -ad+bc-bc+ad&=0,\\
 bc-ad+ad-bc&=0,&
 bd+ac-bd-ac&=0,&
 cd-dc-ab+ab&=0.
\end{align*}
Expansion of the squared norm therefore has coefficient $\Nrm{x}$ on each
of $e^2,f^2,g^2,h^2$ and zero on every mixed term. This proves the first
identity. Substitution of $(a,-b,-c,-d)$ into~\eqref{eq:product} gives the
second, and changing signs leaves squares unchanged, giving the third.

The displayed matrix is only a compact arrangement of sixteen integer
expressions. Each calculation reduces to associativity, commutativity, and
distributivity in coded integer arithmetic; it invokes no linear-algebraic
existence theorem. Those ring laws themselves follow from elementary
quantifier-free arithmetic induction. In ProveIt the signed-pair identity
can be produced by \code{ZT.derive} and \code{ipra\_ring}; see
Section~\ref{sec:formalization}.
\end{proof}

The key consequence is constructive multiplicative closure: given actual
tuples of norms $r$ and $s$, formula~\eqref{eq:product} supplies an actual
tuple of norm $rs$. Taking absolute values gives nonnegative coordinates
without changing the norm.
